\chapter{Antecedentes y estado del arte}
\label{cap:estado_arte}

El presente capítulo recoge los antecedentes técnicos que sustentan el trabajo:
interoperabilidad de información investigadora, formatos de serialización,
mecanismos de validación y técnicas de modelado conceptual. Su finalidad es
comparar alternativas y justificar las decisiones tecnológicas que se utilizarán en
los capítulos posteriores, sin desarrollar todavía el análisis específico del
dominio CVN. Ese análisis se reserva para el capítulo siguiente,
donde se estudiarán sus artefactos oficiales, sus limitaciones y la propuesta Open
CVN construida a partir de ellos.

\section{Representación de información curricular compleja}
\label{sec:estado_arte_representacion_curricular}

La información curricular académica e investigadora constituye un dominio de datos
heterogéneo, con entidades, relaciones, restricciones y niveles de detalle que no
encajan bien en una representación plana. Desde una perspectiva computacional,
esta circunstancia obliga a diferenciar el modelo conceptual del dominio, la
serialización utilizada para intercambiar datos y las herramientas que permiten
procesarlos. Esta separación será una idea recurrente en el resto del documento,
porque permite razonar sobre el significado de los datos sin confundirlo con el
formato concreto en el que se almacenan o transmiten.

En España, CVN proporciona una referencia normalizada para la presentación de
información curricular de investigadores y favorece la interoperabilidad con bases
de datos institucionales~\cite{cvn_web}. En este capítulo se menciona únicamente
como punto de conexión con el dominio de aplicación. La descripción de sus
fuentes, de su representación técnica y de sus limitaciones se reserva para el
capítulo siguiente, con el fin de no anticipar el análisis que corresponde a la
propuesta de solución.

\section{Interoperabilidad curricular y sistemas de información de investigación}
\label{sec:estado_arte_interoperabilidad}

La gestión de información curricular no es un problema exclusivo de CVN. En el
ámbito académico internacional existen diferentes iniciativas orientadas a mejorar
la interoperabilidad entre sistemas de información de investigación. Estos
sistemas, conocidos habitualmente como \textit{Current Research Information
Systems} (CRIS), integran datos sobre investigadores, organizaciones, proyectos,
publicaciones, financiación, resultados científicos y otros elementos relevantes
para la actividad investigadora.

La interoperabilidad en este dominio no depende solo de compartir documentos, sino
también de poder identificar de forma estable a las personas, organizaciones,
resultados y relaciones descritas. ORCID, por ejemplo, proporciona identificadores
persistentes para investigadores y mecanismos de conexión entre dichos
identificadores, contribuciones y afiliaciones~\cite{orcid_about}. En una línea
complementaria, RO-Crate propone un enfoque ligero para empaquetar datos de
investigación junto con sus metadatos~\cite{ro_crate}. Estas iniciativas no
resuelven por sí mismas la representación completa de un currículo CVN, pero
muestran la importancia de combinar datos estructurados, metadatos explícitos e
identificadores reutilizables.

Uno de los modelos más representativos en este ámbito es CERIF, mantenido por
euroCRIS. CERIF se presenta como un formato europeo para la gestión, intercambio y
almacenamiento de información de investigación, con especial atención a la
interoperabilidad entre sistemas CRIS~\cite{eurocris_cerif}. Su principal ventaja
reside en la riqueza de su modelo, que permite representar entidades, relaciones y
dimensiones temporales con un alto grado de precisión. Esta expresividad resulta
adecuada para infraestructuras institucionales complejas, aunque también introduce
un nivel de abstracción y complejidad que puede ser excesivo cuando el objetivo es
definir una representación ligera y directamente procesable de currículos.

Otra aproximación relevante es VIVO, una plataforma orientada a conectar,
compartir y descubrir información sobre investigación mediante tecnologías de Web
Semántica~\cite{vivo_project}. Frente a modelos más tabulares o documentales,
VIVO se apoya en grafos y ontologías para describir investigadores, publicaciones,
organizaciones y relaciones académicas. Esta orientación facilita la publicación
de datos enlazados y la integración con vocabularios semánticos, pero exige una
infraestructura conceptual y tecnológica más compleja que la requerida por una
herramienta local de gestión curricular.

En el contexto español destaca también el proyecto HERCULES y, dentro de él, la
Research Ontology for Hercules (ROH). La red de ontologías HERCULES persigue la
representación semántica de información de investigación en el sistema
universitario español~\cite{hercules_ontologias,roh_documentation}. Además, parte
de su ecosistema se ha relacionado con CVN mediante código y artefactos públicos
orientados a la importación de información curricular~\cite{roh_github,cvn_github}.
Esta proximidad al dominio español convierte a ROH en una referencia relevante,
aunque su orientación ontológica y su dependencia de tecnologías de Web Semántica
hacen que no sea la opción más sencilla para el alcance de este trabajo.

La Tabla~\ref{tab:estado_arte_modelos_interoperabilidad} resume estos tres
enfoques atendiendo al tipo de representación empleada, su principal fortaleza y
la limitación que presenta cada uno frente al objetivo de Open CVN.

\begin{table}[htbp]
  \centering
  \footnotesize
  \renewcommand{\arraystretch}{1.55}
  \caption{Comparativa de modelos de interoperabilidad curricular relacionados con Open CVN.}
  \label{tab:estado_arte_modelos_interoperabilidad}
  \begin{tabular}{p{0.08\textwidth} p{0.26\textwidth} p{0.26\textwidth} p{0.26\textwidth}}
    \hline
    \textbf{Modelo} & \textbf{Enfoque de representación} & \textbf{Fortaleza principal} & \textbf{Limitación frente a Open CVN} \\
    \hline
    CERIF & Modelo europeo orientado a entidades y relaciones para sistemas CRIS & Alta expresividad para relaciones institucionales y dimensiones temporales & Complejidad y coste de adopción elevados para una representación ligera \\
    VIVO & Grafos y ontologías de Web Semántica & Publicación de datos enlazados e integración con vocabularios semánticos & Requiere infraestructura RDF/ontológica no necesaria para el alcance del TFG \\
    ROH & Ontología del ámbito universitario e investigador español (proyecto HERCULES) & Proximidad al dominio CVN español y artefactos de importación ya existentes & Orientación ontológica y dependencia de tecnologías de Web Semántica \\
    \hline
  \end{tabular}
\end{table}

En conjunto, CERIF, VIVO y ROH muestran que la información investigadora puede
modelarse con niveles elevados de riqueza semántica. No obstante, también ponen
de manifiesto una tensión habitual entre expresividad e implementabilidad. Cuanto
más completo y semánticamente rico es un modelo, mayor suele ser el coste de
adopción, mantenimiento, consulta y validación. Para Open CVN, esta observación
resulta determinante: el objetivo no es competir con infraestructuras CRIS ni con
ontologías de investigación completas, sino definir una representación curricular
abierta, validable y suficientemente expresiva para el ámbito del TFG.

\section{Formatos de serialización y validación}
\label{sec:estado_arte_serializacion}

La elección de un formato de serialización condiciona la forma en que los datos se
almacenan, intercambian y procesan. Sin embargo, un formato de serialización no
equivale por sí mismo a un modelo conceptual. Un mismo dominio puede representarse
en XML, JSON, YAML u otros formatos, siempre que exista una definición clara de
las entidades, atributos, relaciones y restricciones que deben preservarse. Por
ello, en este trabajo los formatos se analizan como mecanismos técnicos, no como
la única fuente de significado del sistema.

\subsection{XML y XSD}
\label{subsec:estado_arte_xml_xsd}

Extensible Markup Language (XML) es un lenguaje de marcado definido por el World
Wide Web Consortium y ampliamente utilizado para representar documentos y datos
estructurados. La recomendación XML 1.0 lo define como un perfil de SGML diseñado
para facilitar el intercambio y procesamiento de documentos en la Web, con
atención explícita a la interoperabilidad y a la facilidad de implementación de
procesadores~\cite{w3c_xml_10}. XSD se define, a su vez, como un lenguaje para
describir la estructura y restringir el contenido de documentos XML, aunque la
propia especificación reconoce que algunas aplicaciones pueden necesitar
validaciones adicionales no expresables únicamente mediante el
esquema~\cite{w3c_xsd_11}.

La principal ventaja de XML en este contexto es su capacidad para representar
estructuras jerárquicas complejas y asociarlas a esquemas de validación robustos.
Esto permite comprobar si un documento respeta una estructura determinada y
facilita la interoperabilidad entre sistemas que comparten la misma definición.
Además, el uso de XSD proporciona una fuente formal a partir de la cual pueden
generarse artefactos estructurales, como clases de datos o bindings de acceso
programático. No obstante, XML también presenta limitaciones cuando se utiliza como base directa
de una herramienta de gestión. Su sintaxis es más verbosa que la de formatos más
ligeros, el procesamiento suele requerir bibliotecas específicas y la estructura
del documento no siempre coincide con el modelo conceptual más adecuado para una
aplicación. Por tanto, XML puede ser una fuente técnica fundamental en sistemas de
intercambio estructurado, pero no constituye por sí solo el modelo completo del
dominio representado.

\subsection{JSON y JSON Schema}
\label{subsec:estado_arte_json_schema}

JavaScript Object Notation (JSON) es un formato ligero de intercambio de datos
basado en estructuras de pares clave-valor y listas ordenadas~\cite{json_org_es}.
Su sencillez sintáctica, su integración directa con lenguajes de programación y
su uso extendido en servicios web lo convierten en una alternativa adecuada para
representaciones orientadas al procesamiento automático. En Python, por ejemplo,
un documento JSON puede mapearse de forma natural a diccionarios, listas y tipos
primitivos, lo que simplifica su manipulación en aplicaciones.

La principal dificultad de JSON es que, por sí solo, no define restricciones
estructurales ni semánticas. Para resolver esta limitación se utiliza JSON Schema,
un vocabulario que permite describir la estructura esperada de documentos JSON,
incluyendo tipos, campos obligatorios, restricciones de valores y composición de
objetos. La documentación de JSON Schema lo presenta como un lenguaje declarativo
para definir estructura y restricciones de datos JSON, cuya validación requiere
comprobar una instancia contra un esquema mediante herramientas específicas
~\cite{json_schema_overview}. De este modo, JSON puede combinarse con un mecanismo
de validación formal similar en propósito a XSD, aunque con una sintaxis más
cercana a los formatos habituales de intercambio de datos en aplicaciones
modernas.

En el ecosistema Python, Pydantic complementa esta aproximación al permitir definir
modelos mediante anotaciones de tipo, validar datos de entrada, serializar
estructuras y emitir JSON Schema a partir de los modelos~\cite{pydantic_docs}.
Esta combinación resulta especialmente adecuada para un proyecto que necesita
mantener correspondencia entre un formato de intercambio, reglas de validación y
objetos manipulables desde código. Al mismo tiempo, obliga a distinguir entre la
validación sintáctica o estructural, que puede comprobarse de forma automática, y
las reglas semánticas del dominio, que requieren análisis adicional de las fuentes
CVN.

Para el contexto de Open CVN, JSON resulta especialmente adecuado como formato
final de intercambio y almacenamiento porque permite construir documentos
legibles, validables y compatibles con herramientas habituales de desarrollo.
Esta elección no implica renunciar a una referencia normativa del dominio. Al
contrario, el objetivo consiste en derivar una representación JSON trazable a
partir de fuentes oficiales, preservando la relación entre los datos, sus reglas y
su interpretación curricular.

\subsection{JSON-LD y YAML}
\label{subsec:estado_arte_jsonld_yaml}

JSON for Linking Data (JSON-LD) extiende JSON para facilitar la representación de
datos enlazados mediante contextos semánticos. La recomendación JSON-LD 1.1 lo
define como una serialización basada en JSON para \textit{Linked Data}, pensada
para integrarse con sistemas que ya utilizan JSON y facilitar un camino de
evolución hacia datos enlazados~\cite{w3c_json_ld_11}. Su principal aportación
consiste en permitir que un documento JSON pueda interpretarse también como parte
de un grafo de conocimiento, asociando claves del documento con identificadores
semánticos. Esta característica lo convierte en una opción interesante cuando se
busca compatibilidad con tecnologías RDF u ontologías. Sin embargo, la
incorporación de semántica enlazada también aumenta la complejidad
del sistema. Para un proyecto cuyo objetivo principal es definir una
representación curricular abierta, validable y manejable mediante herramientas
Python, JSON-LD puede resultar más expresivo de lo necesario en una primera
aproximación. Su interés se mantiene como posible línea de evolución futura, en
especial si Open CVN necesitara integrarse con ontologías como ROH o con
infraestructuras de datos enlazados.

YAML, por su parte, es un formato orientado a la legibilidad humana y empleado con
frecuencia en configuración y definición de esquemas~\cite{yaml_org}. Su sintaxis
compacta facilita la lectura manual, aunque puede introducir ambigüedades y
requiere parsers específicos. En el contexto de este trabajo, YAML resulta
relevante principalmente por su uso en herramientas de modelado como LinkML, pero
no se adopta como formato final de intercambio curricular.

La Tabla~\ref{tab:estado_arte_formatos_serializacion} recoge, a modo de síntesis,
la función principal, el mecanismo de validación asociado y el papel que cada uno
de estos formatos desempeña dentro de Open CVN.

\begin{table}[htbp]
  \centering
  \footnotesize
  \renewcommand{\arraystretch}{1.55}
  \caption{Comparativa de formatos de serialización y validación relevantes para Open CVN.}
  \label{tab:estado_arte_formatos_serializacion}
  \begin{tabular}{p{0.13\textwidth} p{0.24\textwidth} p{0.22\textwidth} p{0.24\textwidth}}
    \hline
    \textbf{Formato} & \textbf{Función principal} & \textbf{Validación asociada} & \textbf{Papel en Open CVN} \\
    \hline
    XML + XSD & Representación jerárquica de documentos y datos estructurados & Esquema XSD & Fuente estructural oficial de CVN; no se adopta como formato final de intercambio \\
    JSON + JSON Schema & Intercambio ligero de datos orientado a aplicaciones & JSON Schema, con apoyo de Pydantic en Python & Formato final de intercambio, almacenamiento y validación de Open CVN \\
    JSON-LD & Extensión de JSON para datos enlazados semánticamente & Contextos JSON-LD y vocabularios semánticos & Línea de evolución futura; no se adopta en la versión actual \\
    YAML & Formato legible orientado a configuración y definición de esquemas & Depende de la herramienta que lo consume (por ejemplo, LinkML) & Referencia en herramientas de modelado; no se emplea como formato de intercambio final \\
    \hline
  \end{tabular}
\end{table}

\section{Modelado conceptual y generación de artefactos}
\label{sec:estado_arte_modelado}

La separación entre modelo conceptual y formato de serialización conduce a la
necesidad de utilizar técnicas de modelado que permitan describir el dominio de
forma independiente de una tecnología concreta. En este sentido, Unified Modeling
Language (UML) constituye una referencia ampliamente utilizada para especificar,
visualizar y documentar sistemas mediante clases, atributos, relaciones,
jerarquías y multiplicidades~\cite{omg_uml}. Su utilidad en este trabajo reside
en la posibilidad de representar conceptos curriculares y sus relaciones sin
depender inicialmente de XML, JSON u otro formato de intercambio.

Por este motivo, UML se utilizará como notación principal para representar el
esquema conceptual derivado del análisis del dominio. En particular, los diagramas
de clases permiten expresar entidades curriculares, atributos, asociaciones y
multiplicidades en un nivel intermedio entre la descripción normativa de CVN y los
artefactos técnicos generados posteriormente. Esta representación no sustituye a
los modelos ejecutables ni a los esquemas de validación, pero proporciona una vista
comprensible y tecnológicamente independiente del esquema que se desea obtener.

Cuando el modelo requiere restricciones más precisas que las expresables mediante
clases y relaciones, Object Constraint Language (OCL) permite definir condiciones
formales sobre modelos UML~\cite{omg_ocl}. Estas restricciones pueden utilizarse
para especificar reglas de validez, cardinalidad o consistencia que no siempre
quedan reflejadas de forma natural en un diagrama estructural. Aunque este trabajo
no se basa exclusivamente en UML/OCL como fuente de generación final, estos
conceptos resultan relevantes para justificar la separación entre estructura,
semántica y validación.

En esta línea, existen herramientas que permiten generar artefactos técnicos a
partir de modelos. BESSER, por ejemplo, proporciona un generador de modelos
Pydantic a partir de especificaciones de modelado~\cite{besser_pydantic}. Este
tipo de enfoque muestra una tendencia importante: definir el dominio en un nivel
conceptual o declarativo y obtener después clases, validadores o esquemas de forma
automática. La arquitectura de Open CVN adopta esta idea general de generación y
trazabilidad, aunque la adapta a las fuentes disponibles en el ecosistema CVN y a
las necesidades concretas del proyecto.

Otra herramienta relacionada es LinkML, un lenguaje de modelado orientado a datos
estructurados que permite definir esquemas de forma declarativa y generar
artefactos en distintos formatos~\cite{linkml_io}. LinkML resulta especialmente
interesante por su capacidad para actuar como puente entre modelos, esquemas JSON,
documentación y tecnologías de datos enlazados. No obstante, para el alcance de
este TFG, su adopción completa introduciría una capa adicional de complejidad que
no resulta imprescindible. Por ello, se considera como referencia conceptual, pero
no como dependencia central de la solución desarrollada.

\section{Síntesis del estado del arte}
\label{sec:estado_arte_sintesis}

El análisis realizado permite extraer varias conclusiones relevantes para el
desarrollo de una representación curricular abierta. En primer lugar, los modelos
relacionados con información de investigación, como CERIF, VIVO o ROH, muestran
soluciones expresivas y semánticamente ricas, aunque su complejidad puede resultar
desproporcionada para una herramienta abierta y ligera de gestión curricular. En
segundo lugar, los formatos de serialización deben entenderse como mecanismos de
intercambio y almacenamiento, no como sustitutos del modelo conceptual. En tercer
lugar, UML ofrece una notación adecuada para representar el esquema conceptual de
forma comprensible antes de materializarlo en modelos ejecutables o esquemas de
validación.

Estas observaciones justifican la necesidad de una arquitectura propia que separe
modelo conceptual, serialización y validación, preserve trazabilidad hacia las
fuentes oficiales del dominio y genere una representación final más adecuada para
validación, almacenamiento y exportación. De este modo, el estado del arte no
conduce a adoptar sin modificaciones una ontología completa ni a reproducir
directamente una estructura de intercambio existente, sino a proponer una solución
intermedia: suficientemente formal para ser validable, suficientemente trazable
para conservar el significado curricular y suficientemente ligera para servir como
base de herramientas abiertas.
